% VI/40 reader bridge -- additive mathematical exposition.
\paragraph{Local equations and the Cartier condition.}
Let $X$ be an integral scheme with function field $K(X)$.
A Cartier divisor may be described by an open covering
$X=\bigcup_iU_i$ and nonzero rational functions
$f_i\in K(X)^\times$ such that $f_i/f_j$ is a unit of
$\mathcal O_X$ on every overlap $U_i\cap U_j$.
Two systems describe the same Cartier divisor when they differ
locally by units.  An effective Cartier divisor has local equations
$f_i\in\mathcal O_X(U_i)$ that are nonzerodivisors.
The key requirement is \emph{local principality}: although no
single global equation need exist, the defining ideal must be
generated by one nonzerodivisor on suitable neighborhoods.

\paragraph{Worked example 1: one point on $\mathbb P^1$.}
Let $P=[0:1]\in\mathbb P_k^1$ and take the two standard opens
$U_Y=\{Y\ne0\}$ and $U_X=\{X\ne0\}$.
On $U_Y$, the coordinate $t=X/Y$ vanishes simply at $P$;
on $U_X$, the point $P$ does not occur.  Thus the Cartier
representatives for the effective divisor $D=[P]$ can be chosen as
\[
 f_Y=t\quad\text{on }U_Y,\qquad
 f_X=1\quad\text{on }U_X.
\]
On $U_Y\cap U_X$ we have $t\ne0$, so $f_Y/f_X=t$ is a unit.
The pair glues even though $D$ has no globally defined nonconstant
regular equation on the whole projective line.
The ideal sheaf $\mathcal O_X(-D)$ is generated by $t$ on $U_Y$
and by $1$ on $U_X$.  The corresponding invertible sheaf
$\mathcal O_X(D)$ has local rational generators $t^{-1}$ and $1$.
On the overlap, their transition is multiplication by the unit
$t^{-1}$.  Hence Cartier divisors naturally produce line bundles.

\paragraph{Worked example 2: sections allowed to have a pole.}
By definition, on an integral scheme
\[
 \mathcal O_X(D)(U)
  =\{g\in K(X):\operatorname{div}(g)|_U+D|_U\geq0\}
\]
when the valuation interpretation is available, together with the
zero section.  For $D=[P]$ above, the rational functions $1$ and
$1/t$ are global sections of $\mathcal O_{\mathbb P^1}(D)$:
$1/t$ has precisely one allowed simple pole at $P$ and is regular
at infinity (where $1/t=s$).  They span the two-dimensional
$k$-vector space $H^0(\mathbb P_k^1,\mathcal O(D))$.
Compare this with $H^0(\mathbb P_k^1,\mathcal O)=k$;
the new sections are possible because the divisor has relaxed the
pole condition at a specific point.

\paragraph{Worked example 3: adding two projective lines.}
In $\mathbb P_k^2$, let $L_X=V_+(X)$ and $L_Y=V_+(Y)$.
Each is an effective Cartier divisor.  Their sum is locally
represented by the \emph{product} of defining equations, for example
by $xy$ on the chart $Z\ne0$ with $x=X/Z$ and $y=Y/Z$.
Thus the effective divisor $L_X+L_Y$ is cut out scheme-theoretically
by $XY=0$ (in the homogeneous coordinate ring).  The product
records multiplicities: replacing $XY$ by $X^2Y$ gives
$2L_X+L_Y$.  Addition of divisors corresponds to multiplying
local equations, not to intersecting their set-theoretic supports.

\paragraph{Worked example 4: a Weil divisor that is not Cartier.}
Assume $k$ is algebraically closed of characteristic different from
$2$, and put
\[
 A=k[u,v,w]/(uv-w^2),\qquad X=\operatorname{Spec}A.
\]
This is a normal integral surface with a singular vertex
$\mathfrak m=(u,v,w)$.  The ideal $\mathfrak p=(u,w)$ is a
height-one prime, because $A/\mathfrak p\simeq k[v]$;
therefore it determines a prime Weil divisor $D$.
If $D$ were Cartier near the vertex, its ideal would be generated
by one element in the local ring $A_{\mathfrak m}$.
However, the images of $u$ and $w$ form a two-dimensional
$k$-vector space
\[
 \mathfrak p/\mathfrak m\mathfrak p,
 \qquad \dim_k(\mathfrak p/\mathfrak m\mathfrak p)=2.
\]
Indeed $u,w$ have independent linear terms, while the generators
of $\mathfrak m\mathfrak p$ have total degree at least two.
By Nakayama's lemma $\mathfrak pA_{\mathfrak m}$ needs two
generators.  It is not principal, and $D$ is not Cartier at the
vertex.  A codimension-one cycle need not have a local equation
on a singular scheme.

\paragraph{Comparison with the normal and regular settings.}
On a Noetherian normal integral scheme, a Cartier divisor has an
associated Weil divisor obtained by taking valuation orders of its
local equations.  On a Noetherian regular integral scheme every
height-one prime is locally principal: regular local rings are UFDs.
Consequently every Weil divisor is Cartier on such a scheme.
The cone above shows why replacing ``regular'' with merely ``normal''
is unjustified.  VI/41 measures the remaining obstruction through
linear equivalence and the divisor class group.
